\documentclass{article}
\usepackage{amsmath,amssymb,amsfonts,amsthm}
\usepackage{algorithm}
\usepackage{algpseudocode}
\usepackage{color}

\newtheorem{theorem}{Theorem}
\newtheorem{lemma}{Lemma}
\newtheorem{proposition}{Proposition}
\newtheorem{corollary}{Corollary}

\begin{document}

1. align --- Align multiple equations

Use it for aligned equations, fractions, cases, matrices, and integrals with limits.

\begin{align}
    x^2 -1 = 0\\
    x^3 - 1 = 0
\end{align}

2. gather --- centre multiple equations

each equation centered independently rather than aligned at the equals sign.

\begin{gather*}
    a^2+b^2=c^2\\
    x^2+y^2=z^2
\end{gather*}

3. equation --- a single numbered equation

\begin{equation}
a^2+b^2=c^2
\end{equation}

4. dfrac --- display style fractions

\[
\frac{1}{x+1}
\qquad
\dfrac{1}{x+1}
\]

5. text --- Ordinary words inside mathematics

\[
f(x)=
\begin{cases}
x^2 & \text{if } x\geq 0\\
-x  & \text{if } x<0
\end{cases}
\]

6. cases --- Piecewise functions

\[
f(x)= 
\begin{cases}
    x+1 & \text{if } x > 0\\
    0 & \text{if } x = 0\\
    x-1 & \text{if } x < 0
\end{cases}
\]

7. tag --- Custom equation numbers

\begin{equation*}
E=mc^2 \tag{1.1}
\end{equation*}

8. innertext 

\begin{align*}
x^2-4 &= 0\\
\intertext{Factorizing, we get}
(x-2)(x+2) &= 0
\end{align*}


\pagebreak

\[
x \in \mathbb{R}, \qquad
n \in \mathbb{N}, \qquad
z \in \mathbb{C}
\]

\[
A \subsetneq B, \qquad
C \nsubseteq B, \qquad
\varnothing \subseteq A, \qquad
\emptyset \subseteq A
\]

For every real number \(x\), its square is non-negative.

\[
\forall x \in \mathbb{R} \, \quad x^2 \ge 0
\]

Conclusion

\[
\therefore \quad x^2 \geq 0
\]

\pagebreak

\[
\mathbb{R}, \qquad \mathfrak{g}
\]

$$ \mathbb{R} - \mathbb{Q}  $$

Fraktur letters

\[
\mathfrak{g}, \qquad
\mathfrak{G}, \qquad
\mathfrak{A}
\]

calligraphic letters

\[
\mathcal{F} = \{f : X \to Y\}
\]

\pagebreak

\begin{theorem}
    The sum of two even integer is even.
\end{theorem}


\begin{proof}
Let the two even integers be $2a$ and $2b$.

Their sum is
\[
2a+2b=2(a+b).
\]

Since $a+b$ is an integer, the sum is even.
\end{proof}

Different mathematical environment

\begin{lemma}
    If n is even, then $n^2$ is even.
\end{lemma}

\begin{corollary}
    If n is even, then $n^2+2$ is even.
\end{corollary}

\pagebreak

\begin{algorithm}
\caption{Find the Larger Number}
\begin{algorithmic}[1]
    \State Read $a$ and $b$
    \If{$a > b$}
        \State Print $a$
    \Else
        \State Print $b$
    \EndIf
\end{algorithmic}
\end{algorithm}

Example -

\begin{algorithm}
\caption{Sum of First $n$ Natural Numbers}
\begin{algorithmic}[1]
    \Require A positive integer $n$
    \Ensure The sum $S$
    \State $S \gets 0$
    \For{$i \gets 1$ to $n$}
        \State $S \gets S+i$
    \EndFor
    \Return $S$
\end{algorithmic}
\end{algorithm}

\pagebreak

This is a normal text.

\textcolor{red}{This text is red.}

\textcolor{blue}{Mathematics}

\[
\textcolor{red}{x^2+y^2=z^2}
\]

\end{document}